% GENERATED FILE. Edit canonical lessons, then run npm run generate:books.
% canonical-source-hash: fnv1a64:4fb6e6201e6cee77
% canonical-lessons: KA-C198-ujju, KA-C198-taru, KA-C198-oyyu, KA-C198-hidi, KA-C198-bidu

\chapter{To scrub, To bring, To carry away, To hold, To let go}
\label{ch:ka-to-scrub-to-bring-to-carry-away-to-hold-to-let-go}
\hlchaptermodality{198}

\begin{chapteropening}
\textbf{By the end of this chapter:} \emph{I can say to scrub, to bring, to carry away, to hold and to let go.}

Complete the last lesson of chapter 198: I can say to scrub, to bring, to carry away, to hold and to let go.
\end{chapteropening}

\section[\texorpdfstring{\kn{ಉಜ್ಜು} (ujju)}{ಉಜ್ಜು (ujju)}]{\kn{ಉಜ್ಜು} (ujju) — to scrub, to rub, to brush (teeth)}

\label{lesson:KA-C198-ujju}



\begin{quote}
Before the new one: say the Kannada for a headmaster, then the Kannada for a secretary.
\end{quote}

\subsection*{You'll want to know: \kn{ಉಜ್ಜು}}
\textbf{\kn{ಉಜ್ಜು}} — \emph{ujju} — \textquotedblleft{}to scrub, to rub, to brush (teeth)\textquotedblright{}.

\begin{grammarlens}[title={What you've built}]
One more word for this chapter.
\end{grammarlens}

\subsection*{Your turn}
\begin{itemize}
\raggedright
  \item \emph{Say it:} \emph{ujju}
  \item \emph{Say it:} \emph{ujju}, once more
  \item \emph{Say it:} say \emph{ujju}
  \item \emph{Recall:} say the Kannada for a headmaster, then the Kannada for a secretary, then say \emph{ujju} again
\end{itemize}

\begin{tcolorbox}[breakable,colback=teal!4,colframe=teal!35!black,title={Before you move on}]
What does \kn{ಉಜ್ಜು} mean? (\textquotedblleft{}To scrub, to rub, to brush (teeth)\textquotedblright{}.) Say it once more.
\end{tcolorbox}
\section[\texorpdfstring{\kn{ತರು} (taru)}{ತರು (taru)}]{\kn{ತರು} (taru) — to bring}

\label{lesson:KA-C198-taru}



\begin{quote}
Before the new one: say the Kannada for a secretary, then the Kannada for to scrub.
\end{quote}

\subsection*{You'll want to know: \kn{ತರು}}
\textbf{\kn{ತರು}} — \emph{taru} — \textquotedblleft{}to bring\textquotedblright{}.

\begin{grammarlens}[title={What you've built}]
One more word for this chapter.
\end{grammarlens}

\subsection*{Your turn}
\begin{itemize}
\raggedright
  \item \emph{Say it:} \emph{taru}
  \item \emph{Say it:} \emph{taru}, once more
  \item \emph{Say it:} say \emph{taru}
  \item \emph{Recall:} say the Kannada for a secretary, then the Kannada for to scrub, then say \emph{taru} again
\end{itemize}

\begin{tcolorbox}[breakable,colback=teal!4,colframe=teal!35!black,title={Before you move on}]
What does \kn{ತರು} mean? (\textquotedblleft{}To bring\textquotedblright{}.) Say it once more.
\end{tcolorbox}
\section[\texorpdfstring{\kn{ಒಯ್ಯು} (oyyu)}{ಒಯ್ಯು (oyyu)}]{\kn{ಒಯ್ಯು} (oyyu) — to carry away, to take (somewhere)}

\label{lesson:KA-C198-oyyu}



\begin{quote}
Before the new one: say the Kannada for to scrub, then the Kannada for to bring.
\end{quote}

\subsection*{You'll want to know: \kn{ಒಯ್ಯು}}
\textbf{\kn{ಒಯ್ಯು}} — \emph{oyyu} — \textquotedblleft{}to carry away, to take (somewhere)\textquotedblright{}.

\begin{grammarlens}[title={What you've built}]
One more word for this chapter.
\end{grammarlens}

\subsection*{Your turn}
\begin{itemize}
\raggedright
  \item \emph{Say it:} \emph{oyyu}
  \item \emph{Say it:} \emph{oyyu}, once more
  \item \emph{Say it:} say \emph{oyyu}
  \item \emph{Recall:} say the Kannada for to scrub, then the Kannada for to bring, then say \emph{oyyu} again
\end{itemize}

\begin{tcolorbox}[breakable,colback=teal!4,colframe=teal!35!black,title={Before you move on}]
What does \kn{ಒಯ್ಯು} mean? (\textquotedblleft{}To carry away, to take (somewhere)\textquotedblright{}.) Say it once more.
\end{tcolorbox}
\section[\texorpdfstring{\kn{ಹಿಡಿ} (hiḍi)}{ಹಿಡಿ (hiḍi)}]{\kn{ಹಿಡಿ} (hiḍi) — to hold, to catch}

\label{lesson:KA-C198-hidi}



\begin{quote}
Before the new one: say the Kannada for to bring, then the Kannada for to carry away.
\end{quote}

\subsection*{You'll want to know: \kn{ಹಿಡಿ}}
\textbf{\kn{ಹಿಡಿ}} — \emph{hiḍi} — \textquotedblleft{}to hold, to catch\textquotedblright{}.

\begin{grammarlens}[title={What you've built}]
One more word for this chapter.
\end{grammarlens}

\subsection*{Your turn}
\begin{itemize}
\raggedright
  \item \emph{Say it:} \emph{hiḍi}
  \item \emph{Say it:} \emph{hiḍi}, once more
  \item \emph{Say it:} say \emph{hiḍi}
  \item \emph{Recall:} say the Kannada for to bring, then the Kannada for to carry away, then say \emph{hiḍi} again
\end{itemize}

\begin{tcolorbox}[breakable,colback=teal!4,colframe=teal!35!black,title={Before you move on}]
What does \kn{ಹಿಡಿ} mean? (\textquotedblleft{}To hold, to catch\textquotedblright{}.) Say it once more.
\end{tcolorbox}
\section[\texorpdfstring{\kn{ಬಿಡು} (biḍu)}{ಬಿಡು (biḍu)}]{\kn{ಬಿಡು} (biḍu) — to let go, to leave}

\label{lesson:KA-C198-bidu}



\begin{quote}
Before the new one: say the Kannada for to carry away, then the Kannada for to hold.
\end{quote}

\subsection*{You'll want to know: \kn{ಬಿಡು}}
\textbf{\kn{ಬಿಡು}} — \emph{biḍu} — \textquotedblleft{}to let go, to leave\textquotedblright{}.

\begin{grammarlens}[title={What you've built}]
One more word for this chapter.
\end{grammarlens}

\subsection*{Your turn}
\begin{itemize}
\raggedright
  \item \emph{Say it:} \emph{biḍu}
  \item \emph{Say it:} \emph{biḍu}, once more
  \item \emph{Say it:} say \emph{biḍu}
  \item \emph{Recall:} say the Kannada for to carry away, then the Kannada for to hold, then say \emph{biḍu} again
\end{itemize}

\begin{tcolorbox}[breakable,colback=teal!4,colframe=teal!35!black,title={Before you move on}]
What does \kn{ಬಿಡು} mean? (\textquotedblleft{}To let go, to leave\textquotedblright{}.) Say it once more.
\end{tcolorbox}
